\section{Artifacts, verification, and scope of the result}
\label{sec:artifacts}

\subsection{What is in the archive}

The package contains this article in LaTeX source and compiled PDF, the
Lean specification and implementation plan, and exact-arithmetic Python
diagnostics. It also records the inspected repository snapshot and
reproduction commands. The proof does not depend on an unrecorded
calculation or a numerical table.

\begin{center}
\begin{tabular}{@{}>{\raggedright\arraybackslash}p{0.40\textwidth}>{\raggedright\arraybackslash}p{0.52\textwidth}@{}}
\toprule
File or directory & Purpose\\
\midrule
\code{brun\_ipra.tex}, \code{sections/} &
Complete human-readable argument and the iPRA induction audit.\\
\code{brun\_ipra.pdf} &
Compiled article.\\
\code{lean/BrunSpecification.lean} &
Uncompiled specification and candidate conditional assembly proofs.\\
\code{lean/FORMALIZATION\_PLAN.md} &
Detailed semantic and object-theory implementation route.\\
\code{lean/SOURCE\_LEDGER.md} &
Exact source paths, commit and blob identifiers, toolchain and status.\\
\code{artifacts/finite\_sieve.py} &
Exact finite calculations of the coefficients, quadratic form,
prime-support grouping and partial sums.\\
\code{tests/test\_finite\_sieve.py} &
Ten unit tests checking the implemented finite identities.\\
\code{artifacts/verification.json} &
Reproducible diagnostic results: 1,730 exact arithmetic assertions.\\
\code{README.md}, \code{Makefile} &
Reading guide and reproduction commands.\\
\bottomrule
\end{tabular}
\end{center}

\subsection{What the diagnostics check}

The program uses Python integers and \code{fractions.Fraction}; it uses no
floating-point comparisons. The checks include:
\begin{itemize}
\item prime-factor reconstruction, Mobius cancellation, and root-residue
counts for squarefree moduli;
\item the counting error, \(\lambda_1=1\), the Mobius transform,
\(Q=1/G\), and the bound \(|\lambda_d|\rho(d)\le d\);
\item the partition of \(H_s^2\) by \(\rad(ab)\), and the individual
prime-support group bounds;
\item the square majorant, total sieve error, and twin-prime inclusion;
\item the harmonic and elementary exponential inequalities on sample
parameters;
\item agreement of exact finite Brun sums with factorial-denominator
integer codes, including the double occurrence of \(5\).
\end{itemize}

All ten unit tests passed. The separate JSON diagnostic records 1,730
assertions. These are finite consistency checks on the transcription and
implementation. The universal statements are proved by the argument in
this article; the finite tests do not certify those universal statements,
and do not constitute a Lean proof.

For example, at \(z=5\),
\[
 \sfD_5=\{1,2,3,5\},\qquad G(5)=\frac{14}{3},
\]
and the coefficients are
\[
 \lambda_1=1,\quad
 \lambda_2=-\frac37,\quad
 \lambda_3=-\frac9{14},\quad
 \lambda_5=-\frac5{14}.
\]
The program evaluates the quadratic form exactly as \(3/14=1/G(5)\).
These modest examples are useful checks of signs and normalization,
where a transcription error could invalidate the sieve.

\subsection{Reproduction}

From the extracted package root:
\begin{lstlisting}
python3 -m unittest discover -s tests -v
python3 artifacts/finite_sieve.py --json artifacts/verification.json
latexmk -pdf -interaction=nonstopmode -halt-on-error brun_ipra.tex
\end{lstlisting}
The Python code needs only the standard library.
The PDF build needs a usual LaTeX installation with the packages listed
in the preamble.
Lean compilation must be performed in a workspace providing the pinned
Lean/Mathlib environment, for example:
\begin{lstlisting}
lake env lean /absolute/path/to/lean/BrunSpecification.lean
\end{lstlisting}
That command is a reproduction instruction, not a record of a build
performed here.

\subsection{Final mathematical and formal status}

The completed result is the finite theorem
\eqref{eq:ipra-target}, with its full human-readable sieve proof,
an explicit primitive recursive Cauchy modulus, and a detailed audit of
the permitted induction. It expresses constructive convergence of the
original pair reciprocal series.

The implementation work still needed for a full machine-checked result is
precisely the finite sieve library and the object-language derivation
described in Section~\ref{sec:lean}. The archive does not present that
work as already completed.
